\chapter{移位寄存器——全部变式详解}

\section{变式1：串并转换（SIPO）}

\begin{detailbox}[完整教学]
\textbf{【功能】}串行输入→4拍后并行读取。UART/SPI接收器的基础。
\textbf{【VHDL】}\texttt{reg <= si \& reg(N-1 downto 1); po <= reg;}
\end{detailbox}


\section{变式2：并串转换（PISO）}

\begin{detailbox}[完整教学]
\textbf{【功能】}load并行数据→逐拍串行输出。UART/SPI发送器的基础。
\textbf{【VHDL】}
\begin{lstlisting}
if load='1' then reg <= data_in;
else reg <= '0' & reg(N-1 downto 1); end if;
so <= reg(0);
\end{lstlisting}
\end{detailbox}


\section{变式3：双向移位}

\begin{detailbox}[完整教学]
\textbf{【VHDL】}
\begin{lstlisting}
if dir='1' then reg <= si_r & reg(N-1 downto 1);
else reg <= reg(N-2 downto 0) & si_l; end if;
\end{lstlisting}
\end{detailbox}


\section{变式4：循环移位}

\begin{detailbox}[完整教学]
\textbf{【功能】}数据不丢失——最低位移回最高位。
\textbf{【VHDL】}\texttt{reg <= reg(0) \& reg(N-1 downto 1);}（右循环）
\textbf{【本质】}这就是环形计数器的基础！
\end{detailbox}


\section{变式5：算术移位（保持符号位）}

\begin{detailbox}[完整教学]
\textbf{【逻辑右移】}高位补0 → \texttt{'0' \& reg(N-1 downto 1)}
\textbf{【算术右移】}高位补符号位 → \texttt{reg(N-1) \& reg(N-1 downto 1)}
\textbf{【应用】}有符号数÷2操作。
\end{detailbox}
